\chapter{Introdução}
\label{cap:introducao}

A Engenharia de Software tem como objetivo o desenvolvimento de sistemas que atendam não apenas aos requisitos funcionais, 
mas também a atributos de qualidade que garantam sua sustentabilidade ao longo do tempo. Entre esses atributos, a manutenibilidade 
destaca-se como um dos mais relevantes, uma vez que grande parte dos custos de um sistema está associada à sua evolução e manutenção 
ao longo do ciclo de vida \cite{pressman2010}. Nesse contexto, a modularidade é amplamente reconhecida como um fator fundamental 
para a melhoria da qualidade estrutural de sistemas de software, pois permite a decomposição do sistema em módulos menores e independentes, 
possibilitando que eles sejam modificados de forma isolada e sem que seja preciso conhecer profundamente a implementação de
outros módulos \cite{parnas1972}.

A modularidade está diretamente relacionada a conceitos como coesão e acoplamento. Sistemas com alta coesão e baixo acoplamento 
tendem a apresentar melhor organização estrutural, facilitando sua compreensão, teste e modificação \cite{yourdon1978}. 
Em sistemas orientados a objetos, essas características são fortemente influenciadas pelas decisões arquiteturais adotadas, 
uma vez que a arquitetura define a organização dos componentes e as relações de dependência entre eles \cite{bass2012,sommerville2011}. 
Dessa forma, a escolha de um estilo arquitetural exerce impacto direto na qualidade estrutural do software.

Dentre os estilos arquiteturais amplamente utilizados, destacam-se a arquitetura em camadas e a arquitetura hexagonal. A arquitetura 
em camadas organiza o sistema em níveis hierárquicos, nos quais cada camada possui responsabilidades específicas, promovendo separação 
de responsabilidades \cite{fowler2003}. Por outro lado, a arquitetura hexagonal, também conhecida como Ports and Adapters, 
propõe o isolamento do núcleo da aplicação em relação às dependências externas, favorecendo maior independência tecnológica e 
testabilidade \cite{cockburn2005}. Embora ambos os estilos sejam amplamente discutidos na literatura e utilizados na prática, 
seus impactos na modularidade nem sempre são avaliados de forma quantitativa.

Nesse cenário, métricas de software surgem como instrumentos fundamentais para a avaliação objetiva da qualidade estrutural.
Métricas como coesão e acoplamento permitem quantificar propriedades internas do código, fornecendo base para análise de diferentes 
soluções \cite{fenton2015}. Em particular, métricas estruturais para sistemas orientados a objetos, como as propostas por 
\citeonline{chidamber1994}, são amplamente utilizadas para avaliar atributos internos do software, tais como acoplamento, 
coesão e complexidade, os quais influenciam diretamente a modularidade e a organização estrutural dos componentes.

Apesar da existência de diversos estudos que aplicam métricas de software para análise de sistemas, ainda há limitações na literatura 
no que diz respeito à comparação empírica entre diferentes estilos arquiteturais. Alguns trabalhos concentram-se em arquiteturas 
específicas, analisando seus impactos na modularidade em cenários controlados. Outros estudos realizam comparações entre estilos 
arquiteturais distintos, porém frequentemente baseados em revisões teóricas, sem a utilização de dados extraídos de sistemas reais.

Diante desse contexto, observa-se uma lacuna na realização de estudos experimentais que investiguem, de forma quantitativa, 
o impacto de diferentes estilos arquiteturais na modularidade de sistemas orientados a objetos, utilizando aplicações reais como 
objeto de análise. Assim, este trabalho propõe avaliar a modularidade de sistemas orientados a objetos por meio da aplicação de 
métricas estruturais, realizando um estudo comparativo entre arquitetura em camadas e arquitetura hexagonal. A análise será conduzida 
a partir de dois sistemas desenvolvidos previamente, permitindo a obtenção de evidências empíricas sobre os efeitos das decisões 
arquiteturais na qualidade estrutural do software.

Este trabalho está organizado da seguinte forma: 
o Capítulo \ref{cap:fundamentacao-teorica} apresenta a fundamentação teórica, abordando conceitos de modularidade, arquitetura de software e métricas estruturais; 
o Capítulo \ref{cap:trabalhos-relacionados}, discorre sobre alguns trabalhos relacionados; 
o Capítulo \ref{cap:metodologia} descreve a metodologia adotada;  
o Capítulo \ref{cap:resultados} discute os resultados obtidos; 
e, por fim, o Capítulo \ref{cap:conclusoes} apresenta as conclusões. 

\section{Objetivos}
\label{sec:objetivos}

\subsection{Objetivo Geral}
\label{sec:objetivo-geral}

Avaliar a modularidade de sistemas orientados a objetos por meio da aplicação de métricas estruturais em diferentes estilos arquiteturais.

\subsection{Objetivos Específicos}
\label{sec:objetivos-especificos}

\begin{itemize}
    \item Avaliar a modularidade dos sistemas analisados com base nas métricas estruturais selecionadas;
    
    \item Comparar a modularidade entre os diferentes estilos arquiteturais analisados;
    
    \item Analisar a relação entre decisões arquiteturais e qualidade estrutural do software.
\end{itemize}